Put $AO=x,OB=y$. Then $OD=x$ and $OC=y$ by the perpendicular-bisector conditions. Let $\alpha=\angle DAB,\beta=\angle ABC$. Triangles $AOD$ and $BOC$ are isosceles with base angles $\alpha$ and $\beta$, so both angles are acute. Their central angles are $\angle AOD=180^\circ-2\alpha$ and $\angle BOC=180^\circ-2\beta$.

Since $OA$ and $OB$ are opposite rays, $\angle AOC=2\beta$ and $\angle BOD=2\alpha$. The cosine rule gives
\[AC^2=x^2+y^2-2xy\cos2\beta,\qquad
BD^2=x^2+y^2-2xy\cos2\alpha.\]
Because $x,y>0$ and cosine is one-to-one on $(0,180^\circ)$, the diagonals are equal exactly when $2\alpha=2\beta$, as required.
\begin{asy}
texpreamble("\usepackage{amsmath,amssymb}");
size(240);
defaultpen(fontsize(9pt));
pair A=(0,0),B=(6,0),O=(2,0),D=4*cos(50*pi/180)*dir(50),C=B+8*cos(70*pi/180)*dir(110);draw(A--B--C--D--cycle);draw(A--C,blue);draw(B--D,blue);draw(O--D,gray+dashed);draw(O--C,gray+dashed);draw(O--midpoint(A--D),red+dashed);draw(O--midpoint(B--C),red+dashed);dot(O);label("$A$",A,(-1,-1));label("$B$",B,(1,-1));label("$C$",C,(1,1));label("$D$",D,(-1,1));label("$O$",O,(0,-1));
\end{asy}
